% Section 8.1, from lectures/L08/appendix/a_testing.md.

\section{Testing the Driver on the Host}
\label{c8:sec:testing}\label{c8:app:a}
\chapref{7} built the \code{Uart}: the private register core and every public method of the
interface of \chapref{6}. Nothing has run it. This chapter supplies the missing half --- a transport
that can be scripted, and the provided host suite that drives the real driver over it --- so that a
byte-order or sequencing bug is caught on your machine in milliseconds, long before a logic analyzer
is involved.

The stub is written by you; the host suite is provided and is the check, exactly as the
\file{*_tb.vhd} testbenches were on the VHDL side. The exercises give the precise build
instructions, and this section explains the ideas they rest on.

\subsection{The scripted stub}
\label{c8:sec:scripted_stub}
To test the driver with no hardware, you build a stub that stands in for the
\code{driver::transport::Interface} seam. It records every byte the driver sends and plays back bytes
you have queued for it to return, so a test drives the \emph{real} driver and then asserts the exact
transactions it produced. This is the C++ counterpart of a self-checking testbench: the stub is the
wire, scripted.

The stub is a concrete \code{driver::transport::Interface}. The class is named
\code{driver::transport::Stub}, it publicly inherits the interface so that a \code{Uart} accepts it
wherever the seam is expected, and it overrides the three seam methods, \code{begin()},
\code{transfer()} and \code{end()}. It is injected into the \code{Uart} through the constructor,
exactly like the real transport of \chapref{9}. It lives in \file{include/driver/transport/stub.hpp},
beside the interface it implements, so that the host demo in \file{source/main.cpp} can drive the
driver with no hardware too.

On the \code{MOSI} side it records what the driver sends. Every byte passed to \code{transfer()} is
appended to a log the test can read back, and that log is the exact sequence of bytes the driver put
on the wire, so a test asserts against it to check the five-byte transaction, the command byte and
the byte order. Counting \code{begin()} and \code{end()} calls is part of the contract the provided suite asserts
on: it checks \code{beginCalls()} in eight of its cases, which confirms each transaction was framed
exactly once, and \code{endCalls()} in two, which confirms the two stay balanced.

On the \code{MISO} side it plays back what the driver reads. The test loads a queue of bytes before
the driver runs, each \code{transfer()} call returns the next byte from that queue, and once the
queue is empty \code{transfer()} returns \code{0x00}. That is how a test makes \code{readReg} observe
a chosen register value: queue the four bytes, most significant first, that the peripheral would
have shifted out. The suite reaches for that often enough that the stub owes it a word-at-a-time
helper, which queues a 32-bit register value most significant byte first, and a reset that clears
the queued script between phases.

The result is a self-checking, hardware-free test. It constructs a \code{Stub}, constructs a
\code{Uart} over it, and scripts the responses; it calls the real driver methods under test; and it
then asserts on the stub's recorded bytes and on the driver's return values. No SPI, no AVR and no
bench are involved, and the whole driver is exercised on the host.

\subsection{What the host tests pin down}
\label{c8:sec:host_tests}
Once the driver and stub exist, the provided host suite checks the behaviour that matters:
\begin{itemize}
  \item \textbf{The transactions themselves.} A register read and a register write must produce the
    exact five-byte sequences, with the right command byte and four data bytes \textbf{most
    significant first}. This is the byte-order check.
  \item \textbf{Configuring} must write the baud divider first and then enable the peripheral.
  \item \textbf{\code{write()}} must read \code{STATUS} and write \code{TX_DATA} only when there is
    room, writing nothing when there is not.
  \item \textbf{\code{read()}, on valid data}, must read \code{STATUS}, then \code{RX_DATA}, then
    write \code{RX_POP}, in that order.
  \item \textbf{\code{read()}, on no data}, must report failure and issue \textbf{no}
    \code{RX_POP}.
  \item \textbf{The error flags} must read back, and clearing must write zero to them.
\end{itemize}
The ``no pop on empty'' and ``bytes most significant first'' checks are the two that catch the
mistakes which would otherwise only surface on the bench.

The suite is all-or-nothing by design: it is guarded on all four headers existing, so it reports
nothing to do until \file{register_map.hpp}, \file{uart.hpp}, \file{stub.hpp} and \file{blocking.hpp}
are all present. The first four cases check the \emph{stub}, before any driver case runs, and they
are the ones to get green first. Every \code{Uart} case reads what the stub was scripted to reply, so
a stub bug surfaces as a driver failure: queue a register value least significant first and the
suite reports \code{Uart.WriteWhenReadyPushesTxData} red, pointing at the \code{readReg()} of
\chapref{7}, which is correct and unhelpful.

\subsection{Blocking on top, not inside}
\label{c8:sec:blocking}
\code{write()} and \code{read()} are non-blocking because that is what makes them deterministic to
test: each one reports whether it did anything, and a test can drive both outcomes from the stub's
script. The convenience of ``send this byte, I'll wait'' is real, though, so it is built \emph{over}
the interface as two free functions that spin on the non-blocking core until it succeeds.

Building them that way rather than into the driver keeps the core deterministic and keeps the
spinning where a caller opts into it. They take \code{Interface&} rather than \code{Uart&}, so they
work over any implementation --- including the \code{driver::uart::Stub} of \chapref{6}, which is
what lets the \code{app::EchoNode} of \chapref{10} be tested with no driver at all underneath it.

\subsection{What's ahead}
The exercises in \secref{c8:sec:exercises} build the scripted \code{Stub}, get the host suite green,
add the blocking helpers, and run the provided demo. \chapref{9} then replaces the stub with
\code{AvrSpi} over the ATmega328P's real SPI, and nothing above the seam changes.
